\SourceLeaf{245}%
\Needspace{6\baselineskip}\begin{center}\emph{[\textup{\textsc{Gerling}} to \textup{\textsc{Gauss}}. Marburg, 26 February 1844.]}\end{center}

\{\mbox{.\ .\ .\ .\ .\ .\ .} I have already used the information about \textsc{Bolyai}'s book to order it. Our library, armed with this detailed title, was able to approach the booksellers directly. At the same time, \textsc{Lobachevsky}'s book was ordered from Berlin, so I hope we shall have them before too long. Let me thank you again most heartily for your kind trouble. The Russian papers are of course inaccessible to me, since I have forgotten even how to read them letter by letter, which I learned as a student. In the meantime I have also looked up the paper in \textsc{Crelle}, vol. 17, and everything in it does indeed depend on the j'ai démontré. --- The Russian steppe seems, then, to be suitable ground for these speculations, for \textsc{Schweikart} (now in Königsberg) conceived his »Astral-Geometrie« while he was in Kharkov. \mbox{.\ .\ .\ .\ .\ .\ .}\}

\par\bigskip\Needspace{6\baselineskip}

\begin{center}\textsc{Gauss} to \textsc{Schumacher}. Göttingen, 28 November 1846.\end{center}

\mbox{.\ .\ .\ .\ .\ .\ .\ .\ .} I recently had occasion to look again at \textsc{Lobachevsky}'s short work (Geometrische Untersuchungen zur Theorie der Parallellinien. Berlin 1840, from G. Fincke. 4 printing sheets long). It contains the basic features of the geometry that would have to hold and could hold with strict consistency if Euclidean geometry is not the true one. A certain \textsc{Schweikart}\HistoricalNote{*)}{*)}{Formerly in Marburg, now professor of jurisprudence in Königsberg.} called such a geometry astral geometry, and \textsc{Lobachevsky} called it imaginary geometry. You know that I have held the same conviction for 54 years (since 1792), with a certain later extension that I do not wish to mention here. I have thus found nothing materially new to me
